%!TEX program = uplatex
\RequirePackage{plautopatch}

\documentclass[uplatex,dvipdfmx]{jlreq}

% 余白調整（1ページ収め）
\usepackage[top=18mm,bottom=20mm,left=20mm,right=20mm]{geometry}
\usepackage{amsmath,amssymb}
\usepackage{url}

\pagestyle{empty}
\renewcommand{\baselinestretch}{0.96}

%--------------------------------
% 講演マクロ
%--------------------------------
\newcommand{\talk}[3]{%
\noindent
#1\quad #2\par
\hspace*{2em}#3\par\smallskip
}

%--------------------------------
% タイトル情報
%--------------------------------
\newcommand{\EwMnumber}{第41回}
\newcommand{\EwMtitle}{Euler 生誕300年\\[2mm]\large Euler 数と Euler 類を巡って}
\newcommand{\EwMdate}{2007年9月28日（金）14:30～17:40 ～ 9月29日（土）10:30～17:00}
\newcommand{\EwMplace}{東京都文京区春日1-13-27 中央大学理工学部3号館3階 中央大学高等学校小ホール}

%--------------------------------
% タイトル表示
%--------------------------------
\newcommand{\EwMtitleblock}{
\vspace*{-5mm}
\begin{center}
{\Large ENCOUNTER with MATHEMATICS}

\vspace{2mm}
{\large 数学との遭遇}

\vspace{4mm}
{\Large \bfseries \EwMnumber\ \EwMtitle}

\vspace{4mm}
\EwMdate

\vspace{2mm}
於：\EwMplace
\end{center}

\vspace{2mm}
}

\begin{document}

\EwMtitleblock

%--------------------------------
% プログラム
%--------------------------------
\section*{プログラム}

\subsection*{9月28日（金）}

\talk
{14:30～16:10}
{群の Euler 数 I}
{秋田 利之 氏（北大・理）}

\talk
{16:30～17:50}
{Euler class of surface groups acting on the plane}
{Danny Calegari 氏（Caltech／東工大・情報理工）}

\subsection*{9月29日（土）}

\talk
{10:30～12:00}
{オイラーとトポロジーの誕生}
{松本 幸夫 氏（学習院大・理）}

\talk
{14:00～15:00}
{群の Euler 数 II}
{秋田 利之 氏（北大・理）}

\talk
{15:30～17:00}
{～オイラー類を超える日は来るのか？～}
{森田 茂之 氏（東大・数理）}

\bigskip

17:10～　懇親会（ワイン・パーティー）

%--------------------------------
% 案内文
%--------------------------------
\bigskip

本集会の第41回を以上のような予定で開催いたします。  
非専門家向けに入門的な講演をお願い致しました。  
多くの方々のご参加をお待ちしております。  
講演者による講演内容へのご案内を別紙に添付いたしますのでご覧下さい。

%--------------------------------
% 連絡先
%--------------------------------
\section*{連絡先}

112-8551 東京都文京区春日1-13-27 中央大学理工学部数学教室  

tel : 03-3817-1745  

ENCOUNTER with MATHEMATICS  
\url{http://www.math.chuo-u.ac.jp/ENCwMATH}  

三松 佳彦 : yoshi@math.chuo-u.ac.jp / 高倉 樹 : takakura@math.chuo-u.ac.jp

\end{document}
